\begin{figure}[t]
\centering
\fbox{\begin{minipage}{0.95\columnwidth}
\centering
\textbf{Composition design} $\rightarrow$ target: passivity in chloride, printability (hot-crack resistance), cost/critical-element budget\\[2pt]
$\downarrow$\\[-2pt]
\textbf{Feedstock} $\rightarrow$ gas-atomized powder; solid, cable-type, or powder-cored wire; chemistry compensation for evaporation\\[2pt]
$\downarrow$\\[-2pt]
\textbf{AM process} $\rightarrow$ LPBF / SEBM / DED / WAAM / binder jetting; energy input, scan strategy, preheat, shielding\\[2pt]
$\downarrow$\\[-2pt]
\textbf{As-built state} $\rightarrow$ phases, segregation, texture, defects (pores, cracks, inclusions), surface finish\\[2pt]
$\downarrow$\\[-2pt]
\textbf{Post-processing} $\rightarrow$ HIP, heat treatment (corrosion-qualified), machining/polishing of wetted surfaces\\[2pt]
$\downarrow$\\[-2pt]
\textbf{Marine performance} $\rightarrow$ passivation, pitting/crevice, SCC and corrosion fatigue, cavitation, biofilm coupling, galvanic behavior\\[2pt]
$\downarrow$\\[-2pt]
\textbf{Qualification} $\rightarrow$ defect acceptance, electrochemical screening, long-term exposure, class approval\\[2pt]
$\downarrow$\\[-2pt]
\textbf{Deployment} $\rightarrow$ seawater systems, heat exchangers, propulsion, repair cladding; monitoring and life-cycle review
\end{minipage}}
\caption{Process--structure--property--qualification chain for AM HEAs in marine service. The framework organizes the review: each stage is a control point, and the weakest link (long-term seawater evidence and standards coverage) defines the deployment roadmap.}
\label{fig:heam_framework}
\end{figure}
